\section{Introduction}

\subsection{Regulatory drivers make reach-scale PFAS prediction a decision tool}
The regulatory basis for managing per- and polyfluoroalkyl substances (PFAS) in surface water has shifted decisively since 2024, moving the practical demand placed on environmental scientists from detection toward quantitative, reach-scale prediction of fate, transport, and source contribution. In April 2024 the U.S.\ Environmental Protection Agency (EPA) promulgated the first PFAS National Primary Drinking Water Regulation, establishing legally enforceable Maximum Contaminant Levels (MCLs) of $4.0$~ng~L\textsuperscript{-1} for both perfluorooctanoic acid (PFOA) and perfluorooctane sulfonate (PFOS), $10$~ng~L\textsuperscript{-1} individually for PFHxS, PFNA, and HFPO-DA, and a unitless Hazard Index of one for their mixtures \citep[89~Fed.~Reg.~32532;][]{epa2024mcl}. As of mid-2026 the landscape is in flux but not retracted: in May 2025 EPA announced it would retain the enforceable PFOA and PFOS MCLs while extending the compliance deadline to 2031, and has since proposed to rescind the MCLs for PFHxS, PFNA, HFPO-DA, and the Hazard Index \citep{epa2025mcl}. The two $4.0$~ng~L\textsuperscript{-1} limits---the values most relevant to source-water characterization---therefore remain in force, at thresholds near the analytical detection floor, so the source waters feeding drinking-water intakes must be characterized at trace levels, a watershed-loading question rather than a treatment-plant question alone. In the same period EPA designated PFOA and PFOS as hazardous substances under CERCLA \citep[89~Fed.~Reg.~39124;][]{epa2024cercla}, whose strict, joint-and-several liability makes source apportionment---which release reached which receptor, by which pathway, and in what proportion---a legal determinant of cost recovery rather than only a scientific question. Finally, under Clean Water Act \S303(d), a PFAS impairment listing triggers a Total Maximum Daily Load, defined as the sum of point-source wasteload allocations, nonpoint load allocations, and a margin of safety \citep{epatmdl}; PFAS-specific TMDLs remain emergent rather than routine, in part because the wasteload/load split requires a transport tool that can resolve the nonpoint contributions monitoring cannot isolate. State in-stream criteria already make the demand concrete: Michigan's Rule~57 surface-water values for water bodies protected as drinking-water sources are $12$~ppt for PFOS and $66$~ppt for PFOA \citep{egle2022rule57,egle2023rule57}---the same quantity a reach-resolved model predicts.

\subsection{Why monitoring and machine learning are necessary but not sufficient}
These drivers expose capabilities that monitoring and statistical occurrence models, by construction, cannot supply, and the gap is one of capability rather than goodness-of-fit. A grab sample integrates surface runoff, lateral flow, groundwater discharge, and sediment-bound transport, each with distinct spatial origins, so only a model that separately routes each pathway can disaggregate an observed concentration back to its sources \citep{rafiei2023}. PFAS monitoring is, moreover, structurally sparse and analyte-limited: a synthesis of more than 45{,}000 global samples shows that conventional targeted suites systematically underestimate total PFAS burden \citep{grunfeld2024}, and catchment-scale export cannot be constrained without flow-weighted mass-loading data that grab sampling does not provide \citep{byrne2024}. The dominant continental-scale response has been statistical: machine-learning models predict where PFAS are likely to occur from landscape and proximity-to-source covariates---at national scale for groundwater \citep{tokranov2024}, for drinking-water sources \citep{fernandez2023}, and, exploiting the relational structure among wells, with graph neural networks for unmonitored supply wells \citep{rafiei2026gnn}. These tools are valuable for triage and achieve strong predictive skill where observations are dense, but they are correlative: they predict a probability or class of occurrence, not a mass-conserving concentration in a defined reach, and they cannot represent the source-control or remediation interventions that constitute a regulator's actual levers, because an intervention acts on a \emph{pathway} rather than directly on a concentration. We therefore frame the process model as the causal and forecasting tier that sits atop monitoring and machine learning---complementary to their occurrence-triage strength rather than a replacement---and situate it within the long program on prediction in ungauged basins \citep{hrachowitz2013}, since a parameter-transferable process model is the only route to a defensible concentration on the unmonitored majority of a high-resolution stream network.

\subsection{Prior art and the gap}
Process-based watershed-scale PFAS modeling remains comparatively immature and is dominated by a single research lineage. A pair of reviews catalogued the field and concluded that, although individual process models existed for the vadose zone, groundwater, and receiving waters, no framework coupled these compartments into a continuous watershed-scale mass balance, identifying the land--stream interface and compound-specific partitioning as the principal gaps, and demonstrating in a companion case study the heavy data and calibration burden that integration imposes \citep{raschke2022overview,raschke2022case}. Building on this foundation, a distributed, watershed-scale hydrogeochemical model for PFOS resolved leaching, runoff, lateral flow, and sediment-bound transport from point and nonpoint sources for source attribution \citep{rafiei2023}; to our knowledge the first distributed watershed-scale PFAS surface-water model, it is the most direct antecedent to the present work and follows the established solute-coupling tradition of SWAT--MODFLOW \citep{bailey2016} and multispecies reactive transport \citep{clement1998}. A recent, independent review reaches a consonant conclusion, finding persistent advances-and-gaps in PFAS modeling for watersheds and receiving waters \citep{zhang2025}. Two boundaries of this prior art define our gap. First, process-based PFAS modeling in the SWAT family is bespoke and, where it has matured for point-source organics, has targeted readily degradable pharmaceuticals rather than persistent, sorption-controlled PFAS \citep{salo2025}; the major process-based PFAS modeling investment elsewhere is subsurface- and source-zone-focused, serving site remediation rather than reach-resolved surface-water fate. Second, none of this work is native to the \emph{modern}, free-form SWAT\textsuperscript{+} engine, resolved at the high reach density of national high-resolution hydrography, automatically generated, or validated against the engine's own established constituent mathematics. The PFAS-specific physics that any such model needs is, moreover, well grounded: beyond Freundlich solid-phase sorption, PFAS accumulate at the air--water interface as interfacially active surfactants, a mechanism that can dominate vadose-zone retention \citep{brusseau2018,guo2020}, while the in-stream behavior follows the organic-carbon partitioning of hydrophobic organics \citep{karickhoff1979} and the settling, resuspension, diffusion, and burial of toxic-organic surface-water modeling \citep{chapra1997}. National source inventories now supply the forcing layer that such a model needs, mapping tens of thousands of presumptive PFAS sites \citep{salvatore2022}. What has been missing is a transferable, automated, high-resolution, process-based PFAS surface-water fate-and-transport tool native to the open-source watershed-modeling stack; this is the gap the present work fills.

\subsection{Contribution and scope}
We make three contributions. First, we implement a complete PFAS surface-water fate-and-transport engine in modern SWAT\textsuperscript{+} \citep{bieger2017}: a double-precision soil three-phase equilibrium solver (aqueous, Freundlich solid, Langmuir air--water interface), land-phase mobilization through runoff, lateral flow, leaching, and sediment, a dedicated in-stream routine that is numerically identical to the SWAT\textsuperscript{+} pesticide channel solver on its linear-partition subset, mass-conserving reservoir mixing, per-reach point sources, and daily reach-resolved concentration output with a closing mass balance. Second, we couple the engine to automatically generated, high-resolution SWAT\textsuperscript{+} models built from NHDPlus High Resolution hydrography, so that ambient PFAS observations are assigned to their hydrologically correct reaches by drainage area rather than to coarse aggregated subbasins. Third, we demonstrate the coupled system on the Rogue River watershed (Michigan), home to the Wolverine World Wide tannery and House Street disposal site, and show that a calibrated model reproduces the observed downstream PFOS gradient measured at 31 state ambient-monitoring stations. We scope this first engine version to PFOS, treat the demonstration as a single-watershed test rather than a national validation, and manage---rather than eliminate---equifinality through physically grouped, class-based parameters. The remainder of the paper details the engine and its correctness checks (Section~\ref{sec:engine}), the high-resolution model generation and Rogue study area (Section~\ref{sec:model}), the calibration design (Section~\ref{sec:calib}), the results (Section~\ref{sec:results}), and their interpretation and limits (Sections~\ref{sec:discussion}--\ref{sec:conclusions}).
